% Generated draft: requires mathematical and source review.
\documentclass[11pt]{article}
\usepackage[margin=25mm]{geometry}
\usepackage{fontspec}
\setmainfont{Latin Modern Roman}
\usepackage{amsmath,amssymb,mathtools,unicode-math}
\setmathfont{Latin Modern Math}
\usepackage{enumitem,needspace,xurl,hyperref}
\setlength{\parindent}{0pt}
\setlength{\parskip}{5pt}
\setlength{\emergencystretch}{3em}
\newcommand{\N}{\mathbb{N}}
\newcommand{\Z}{\mathbb{Z}}
\newcommand{\Q}{\mathbb{Q}}
\newcommand{\R}{\mathbb{R}}
\newcommand{\C}{\mathbb{C}}
\newcommand{\F}{\mathbb{F}}
\newcommand{\A}{\mathbb{A}}
\newcommand{\PP}{\mathbb{P}}
\newcommand{\E}{\mathbb{E}}
\newcommand{\eps}{\varepsilon}
\newcommand{\dd}{\,\mathrm{d}}
\begin{document}
\section*{UnsolvedMath ID: 66}
\textbf{Source problem code:} catalogue entry
\subsection{Definitions and mathematical statement}
DARPA's Mathematical Challenge Two (2007), ``The Dynamics of Networks,'' asks for the development of high-dimensional mathematical methods to model and predict the behavior of large-scale distributed networks. The supplied examples include the internet, social networks, and biological networks, and the background emphasizes network dynamics.

This is a research programme, not a specified mathematical proposition. The source does not fix a class of networks, a state space or evolution law, the meaning of ``large-scale'' or ``high-dimensional,'' available observations, a prediction target, or a quantitative accuracy or performance requirement. Consequently it supplies no domain and quantifiers for a theorem and no formal completion criterion. A precise instance would require those choices; none is selected here.

The local registry labels the problem open. Its literature triage dated 2026-08-17 labels the agenda partially solved in the sense of partial progress, cites DARPA's later Models, Dynamics and Learning programme as pursuing rigorous data-driven models for non-equilibrium dynamics, including communication and social systems, and explicitly states that a full solution or refutation is not applicable without a success condition. These are dated source assessments, not a claim that the agenda has a mathematically defined solved status.
\subsection{Short English statement}
Develop mathematics for modelling and predicting the dynamics of massive distributed networks. The supplied challenge is an open-ended research agenda: a network model, prediction task, data regime, and quantitative guarantee must first be specified to obtain a precise mathematical problem.
\subsection{Sources}
\begin{itemize}
\item DARPA DSO, Broad Agency Announcement 07-68, Mathematical Challenge Two (2007), as supplied in the record: \url{https://web.math.utk.edu/~vasili/refs/darpa07.MathChallenges.html}.
\item DARPA, Models, Dynamics and Learning programme page, accessed 2026-08-17 according to the supplied literature triage: \url{https://www.darpa.mil/research/programs/models-dynamics-and-learning}.
\item Local registry \texttt{data/unsolvedmath/problems.json}, permanent catalogue entry, source code \texttt{catalogue entry}, ``The Dynamics of Networks,'' attributed to DARPA (2007); statement, background, and literature triage dated 2026-08-17. Dataset attribution: Ulam AI and the UnsolvedMath Contributors.
\end{itemize}
\end{document}
